\appendix
\section{复算索引与论文使用说明}
论文正文数值由仓库中的四份结果报告及其 CSV/JSON 表核对；程序分别为 \texttt{code/problem1/problem1.py}、\texttt{code/problem2/problem2.py}、\texttt{code/problem3/problem3.py}、\texttt{code/problem4/problem4.py}，输出分别在 \texttt{outputs/problem1} 至 \texttt{outputs/problem4}。每问的独立验收程序及报告位于对应的 \texttt{code/problemN/verify\_problemN.py} 和 \texttt{reports/问题N/验证验收/}。论文引入的 PDF 图均引用这些程序已生成的文件；论文未重新构造实验点或补填缺失评分。第四问另有前沿分位与任务标尺补充程序 \texttt{code/problem4/problem4\_frontier\_sensitivity.py}，其独立回代脚本为 \texttt{code/problem4/verify\_problem4\_frontier\_sensitivity.py}，审查意见核查见 \texttt{reports/问题四/验证验收/IMAGE\_SUGGESTION\_AUDIT.md}。

正文保留三类待核验信息：参赛队号、学校及队员信息不在数据附件中，因此当前匿名稿不填虚构署名；正式提交前按竞赛当届规则处理封面。题目附件中的版权、来源和模型开放许可仍应按最终提交要求再次人工核查。若重跑数据与现有报告不一致，应先更新结果报告，再同步修改论文表格与图注，而非只改文字。

\section{主要符号}
\begin{tabular}{p{2.1cm}p{10.5cm}}\toprule
符号 & 含义\\\midrule
$Q_A$ & A1 指标体系下的相对质量评分\\
$Q_B$ & B6 原生质量水平，本文只在其观测支持内使用\\
$p$ & 17 域占比向量，$\sum_i p_i=1$\\
$N_B,D_B$ & 以十亿为单位的参数量和训练 token 数\\
$L_0,L_1$ & B1 基线标度式及 B6 来源校准后的质量式\\
$C,L_{\rm ctx}$ & FLOPs 预算及外生上下文长度\\
$S$ & C1 同协议六任务平均能力分数\\\bottomrule
\end{tabular}

\section{AI 工具使用披露}
本项目使用 OpenAI Codex 协助梳理题目与数据字段、讨论模型方案、编写和修改分析程序、运行结果核验、绘制图表以及起草本论文文字。论文中的数值和图形须能由题目附件及仓库程序复算；Codex 输出本身不作为数据来源或学术依据。参赛队在正式提交前须逐式、逐表和逐图独立核验，并对模型选择、计算过程与最终论证负责。
